\section{Closure of the classical growth argument}
\label{sec:classical-closure}

We complete the proof of Theorem~\ref{thm:classical-growth}.  Fix
$0<\gamma<1$ and suppose for a contradiction that $\beta>\gamma$.
The least-aspect construction
and isotropic exclusion leave the matched process with $m>0$ used in
Sections~\ref{sec:characteristic-plane} and~\ref{sec:rate-transfer}.
We keep the same saturated homogeneous component, so its functions
$b,S,\Phi$ and its root level $\tau$ are the ones selected before the
trace point and block scales.

On the aperture region \eqref{eq:matched-aperture-region}, the potential
$\Phi$ satisfies the obstacle \eqref{eq:closure-obstacle} and meets it
along $i=b(u)$.  The trace theorem supplies the plane rates
$x,y,w,z,D$, with the single middle rate $q=a_U$ when $m=1$.
Proposition~\ref{prop:rate-transfer} constrains these rates at generic
plane points with $i>0$.  At almost every contact time, $p=-b'(u)$ belongs to
$[0,1+m]$ and $s=S'(u)$ belongs to $(0,1]$.
We now transfer the generic rate constraints to tuples on the inward
sides of the contact curve, obtain a positive tangent rate there, and
combine them with the contact relation.

All information and entropy expressions in this section are limiting
normalized quantities.  In every comparison involving an additional
fixed depth, its finite-scale ambiguity tends to zero before that depth
is sent to zero.

\subsection{Rate tuples at a contact point}

For distinct clocks put
\begin{align}
 B&=x+y+w+z-2,\\
 F(P)&=2P-m-Px+(1-P)y+(m-P)w+(1+m-P)z.
 \label{eq:closure-F}
\end{align}
For tied clocks these definitions mean
\begin{equation}
 B=x+q+z-2,\qquad
 F(P)=2P-1-Px+(1-P)q+(2-P)z.
 \label{eq:closure-F-tied}
\end{equation}
In either case $F'(P)=-B$, with the rates held fixed.

\begin{lemma}[Contact tuples]
\label{lem:closure-contact-tuples}
At almost every $u_0\in(0,1)$, set
\[
 b_0=b(u_0),\qquad p=-b'(u_0),\qquad s=S'(u_0).
\]
The following assertions hold on each side of the contact curve that
lies inward from the aperture boundary.  Minus denotes smaller $i$ and
plus denotes larger $i$.
\begin{enumerate}
\item Rescale area measure on a fixed bounded region of positive area
in a tangent half-plane and form the joint distribution of the generic
plane rates.  Tuples in a subsequential limiting support obey the
capacity bounds and all the implications
of Proposition~\ref{prop:rate-transfer}, with the fixed value $s$.
Every data rate with clock $c_j\ne p$, and also $D$, has a fixed value
on each such support, equal to its one-sided strong trace.  Here
$c_j\in\{0,1,m,1+m\}$, with the two middle clocks grouped when $m=1$.
\item Every tuple on either side satisfies
\begin{equation}
 F(p)+D=2s-\beta-\gamma m.
 \label{eq:closure-contact-identity}
\end{equation}
\item There is a minus tuple with $B\le0$ whenever that side is
admissible, and there is a plus tuple with $B\ge0$ whenever that side
is admissible.
\item If both sides are admissible and $p<1+m$, then
\begin{equation}
 D^- -D^+\ge (1+m-p)(z^+-z^-).
 \label{eq:closure-jump}
\end{equation}
\end{enumerate}
If $b_0=0$, then $p=0$ and the plus side is admissible.  If
$b_0+(1+m)u_0=L$, then $p=1+m$ and the minus side is admissible.
\end{lemma}

\begin{proof}
Choose a common differentiability point of $b,S$ which is a Lebesgue
point of $S'$ and a point of the strong curve traces in
Lemma~\ref{lem:curve-traces}.  The assertions concerning aperture
boundaries follow by differentiating the nonnegative Lipschitz
functions $b(u)$ and $L-b(u)-(1+m)u$ at a zero.  Their derivative is
zero at every such interior differentiability point.

For $h\downarrow0$ take a locally uniform subsequential limit
\[
 \phi_h(\xi,\eta)
 =\frac{\Phi(b_0+h\xi,u_0+h\eta)-\Phi(b_0,u_0)}{h}
 \longrightarrow\phi(\xi,\eta).
\]
Uniform local Lipschitz bounds give this compactness.  On the tangent
line $\xi+p\eta=0$, contact and differentiability give
\begin{equation}
 \phi(-p\eta,\eta)=k\eta,
 \qquad k=-2p+(1-\gamma)m+2s-\beta.
 \label{eq:closure-tangent-contact}
\end{equation}
On each inward half-plane the obstacle becomes
\begin{equation}
 \phi(\xi,\eta)\ge
 \ell(\xi,\eta):=2\xi+((1-\gamma)m+2s-\beta)\eta.
 \label{eq:closure-linear-obstacle}
\end{equation}
These assertions include the boundary line by continuity.  They follow
first on compact subsets of the limiting aperture domain and then on
its boundary by the common Lipschitz bound.

The plane derivative identities \eqref{eq:trace-plane-derivatives} are
\[
 \Phi_i=x+y+w+z,\qquad
 \Phi_u=y+mw+(1+m)z+D,
\]
with the grouped interpretation at $m=1$.  By
Lemma~\ref{lem:curve-traces}, the expression
\[
 \Phi_u-p\Phi_i
 =-px+(1-p)y+(m-p)w+(1+m-p)z+D
\]
converges strongly in mean on each tangent half-ball to a constant.
Indeed every rate appearing with nonzero coefficient is transverse to
the contact curve; a possibly untraced rate has $c_j=p$ and coefficient
zero.  Passing to weak derivatives of $\phi_h$ therefore shows that
$(\partial_\eta-p\partial_\xi)\phi=k_\pm$ is constant on each
half-plane.  In coordinates $r=\xi+p\eta,\eta$, this implies
$\phi(r-p\eta,\eta)=k_\pm\eta+H_\pm(r)$ locally.  Continuity at
$r=0$ and \eqref{eq:closure-tangent-contact} give $k_\pm=k$.
Adding $2p-m$ proves \eqref{eq:closure-contact-identity}.

Here is a precise way to choose tuples and the horizontal signs.  Fix
a bounded strip
\[
 \Omega_+=\{(\xi,\eta):|\eta|<a,\ 0<\xi+p\eta<a\}
\]
on an admissible plus side, or its counterpart $\Omega_-$ with
$-a<\xi+p\eta<0$.  Sample the generic rate vector at
$(b_0+h\xi,u_0+h\eta)$ with uniform area measure on that strip,
including $S'(u_0+h\eta)$ in the vector.  Intersections with the
actual inward domain, if necessary, change a set of relative area
$o(1)$.  Boundedness gives a weakly convergent subsequence of these
probability laws.  The strong traces fix every transverse rate and
$D$ on the support, and the Lebesgue-point property fixes the last
coordinate to $s$.  Each strict-hypothesis implication in
Proposition~\ref{prop:rate-transfer} persists on this support: a
violation with a strictly positive source or a strictly subunit target
would hold on an open neighborhood disjoint from every sufficiently
late sampled law.  Thus all these tuples have the claimed generic
constraints and the identity just proved.

For almost every fixed $\eta$, horizontal integration of
\eqref{eq:closure-linear-obstacle}, using equality at $r=0$, gives
\[
 \int_{-a}^{0}\partial_\xi(\phi-\ell)(r-p\eta,\eta)\,dr\le0,
 \qquad
 \int_{0}^{a}\partial_\xi(\phi-\ell)(r-p\eta,\eta)\,dr\ge0
\]
on the corresponding admissible sides.  Since the limiting horizontal
derivative is the weak limit of $B$, the mean of $B$ in the minus
tuple law is nonpositive and in the plus tuple law is nonnegative.
A continuous function on a compact support cannot have nonpositive
mean while being strictly positive everywhere, nor nonnegative mean
while being strictly negative everywhere.  This supplies the two
required tuples.  The same construction can be made for any fixed
bounded subset of a tangent half-plane; only finitely many such
subsequences will be used below.

Finally, the plane measure inequality \eqref{eq:trace-z-D} is
\[
 (\partial_u-(1+m)\partial_i)z\ge\partial_iD.
\]
When $p<1+m$, both $z$ and $D$ have strong traces.  Rescaling this
inequality at the contact point leaves piecewise constant limits on
the two half-planes.  Their distributional derivatives across
$\xi+p\eta=0$ have coefficients
$(p-1-m)(z^+-z^-)$ and $D^+-D^-$, respectively.  Positivity of the
limiting measure gives \eqref{eq:closure-jump}.
\end{proof}

\subsection{Aperture refinement forces a positive tangent rate}

\begin{lemma}[Positive $X$ rate]
\label{lem:closure-positive-x}
At the contact points of Lemma~\ref{lem:closure-contact-tuples},
$x^-\ge\gamma$ if $p>0$.  If $p=0$, an admissible plus tuple can be
chosen with $x\ge\gamma>0$.
\end{lemma}

\begin{proof}
Fix such a point and let $E_0$ be its contact observation, with labels
$C_{u_0},R_{mu_0}$ and data depths
\[
 b_0,\quad b_0+u_0,\quad b_0+mu_0,\quad b_0+(1+m)u_0.
\]
Keep these labels and refine $X,Y$ by a depth $d$, where
$0<d<mu_0$.  The resulting observation is an admissible aperture with
width depths $(b_0+d,b_0+mu_0)$; the old orientation prefix is an
allowed extra label.  Passing to the required coarser orientation
prefix of depth $mu_0-d$ creates normal velocity uncertainty at most
$O(\eps^{mu_0-d}\eps^{b_0+d})=O(\eps^{b_0+mu_0})$; the position
calculation has the same extra factor $\eps^{u_0}$ on both sides.
Thus the old-frame boxes and a standard aperture test have bounded
ambiguity.  This aperture refinement is available also when $m=1$;
it does not use the grouped auxiliary profile.  Its weight increases
by $(1+\gamma)d$.
Comparing its universal score bound with the saturated score at
$E_0$ yields
\[
 \Ent(X_{b_0+d},Y_{b_0+u_0+d}\mid E_0)\ge(1+\gamma)d.
\]
The $Y$ refinement, even after conditioning on the new $X$, has
entropy at most $d$, by scalar box counting.  Consequently the
nondecreasing Lipschitz profile
\begin{equation}
 g(d):=\Ent(X_{b_0+d}\mid E_0)
 \quad\text{satisfies}\quad g(0)=0,\qquad g(d)\ge\gamma d.
 \label{eq:closure-x-entropy}
\end{equation}
These profiles can be extracted along the same sequence, first on
countably many depths and then by their Lipschitz extension.

Choose $0<\lambda<\min\{1,m\}$, put $v=u_0-h$, and consider
\begin{equation}
 b_0\le i\le b_0+\lambda h,
 \qquad h>0\text{ sufficiently small}.
 \label{eq:closure-x-wedge}
\end{equation}
We claim that $E_0$ together with $X_i$ determines the plane
observation $E(i,v)$ to zero information cost.  The earlier labels
are coarsenings of those in $E_0$.  With
\[
 c=C_v-C_{u_0},\qquad r=R_{mu_0}-R_{mv},
 \qquad |c|\lesssim\eps^v,\quad |r|\lesssim\eps^{mv},
\]
the change of coordinates is
\[
 Y_v=Y_{u_0}+cX,\quad W_v=W_{u_0}+rX,\quad
 Z_v=Z_{u_0}+cW_{u_0}+rY_{u_0}+crX.
\]
Only errors within the recorded boxes matter, since all their centers
and the label shifts are known.  The old $Y$ error fits depth $i+v$
because $i-b_0<h$, and the old $W$ error fits depth $i+mv$ because
$i-b_0<mh$; the terms $cX,rX$ fit those depths using $X_i$.
For $Z_v$ the four error depths are, respectively,
\[
 b_0+(1+m)u_0,\quad b_0+(1+m)u_0-h,\quad
 b_0+(1+m)u_0-mh,\quad i+(1+m)v.
\]
All are at least the target depth $i+(1+m)v$, by
$i-b_0<\min\{1,m\}h$.  Thus the asserted determination holds with
bounded box ambiguity, which vanishes in normalized entropy.

For a further fixed $t>0$ with $i+t<b_0+mu_0$, conditioning on
$E_0,X_i$ therefore gives
\begin{equation}
 \Ent(X_{i+t}\mid E(i,v))
 \ge g(i-b_0+t)-g(i-b_0).
 \label{eq:closure-x-increment}
\end{equation}
The left side is the increase of particle information when just
$L_X$ is raised by $t$ at the plane point.  This is an inward
direction of the resolution cone.  At almost every $(i,v)$ the
affine cone blow-up in Theorem~\ref{thm:characteristic-trace} supplies
a sequence $t\downarrow0$ along which this increase divided by $t$
tends to $x(i,v)$.  At almost every $i$, the right side divided by
$t$ tends to $g'(i-b_0)$.  Hence
\[
 x(i,u_0-h)\ge g'(i-b_0)
\]
at almost every point of \eqref{eq:closure-x-wedge}.  Fubini's theorem
and \eqref{eq:closure-x-entropy} imply, for almost every sufficiently
small $h$,
\begin{equation}
 \int_{b_0}^{b_0+\lambda h}x(i,u_0-h)\,di
 \ge g(\lambda h)-g(0)\ge\gamma\lambda h.
 \label{eq:closure-x-average}
\end{equation}

Integrate also over $h\in[\rho,2\rho]$ and send $\rho\downarrow0$.
If $p>0$, take $\lambda<p$ as well.  Differentiability gives
$b(u_0-h)=b_0+ph+o(h)$, uniformly in this range, so the region lies
on the minus side for small $\rho$.  Since $c_X=0\ne p$, the strong
trace of $x$ makes its average tend to $x^-$.  Inequality
\eqref{eq:closure-x-average} gives $x^-\ge\gamma$.

If $p=0$, the same region is on the plus side apart from area
$o(\rho^2)$.  Removing that part changes the average by $o(1)$,
since $0\le x\le1$.  A subsequential tuple law from the remaining
region therefore has mean $x$ at least $\gamma$, and its compact
support contains a tuple with $x\ge\gamma$.  All its other
transverse components are the same plus traces used in
Lemma~\ref{lem:closure-contact-tuples}.  This proves the assertion
without requiring a trace for the characteristic component $x$.
\end{proof}

\subsection{The finite sign calculation}

\begin{lemma}[Algebraic rate bounds]
\label{lem:closure-sign-algebra}
Let $m>0$, $0<s\le1$, and let a rate tuple obey the capacity bounds
and the implications of Proposition~\ref{prop:rate-transfer}.  Put
$p_*=(1+m)/2$.  Then
\begin{align}
 x>0,\ z<1&\quad\Longrightarrow\quad F(p_*)\ge1+s,
 \label{eq:closure-midpoint}\\
 x>0,\ z=1&\quad\Longrightarrow\quad F(0)\ge1+s,
 \label{eq:closure-left-end}\\
 x=0,\ z<1&\quad\Longrightarrow\quad F(1+m)\ge1+s.
 \label{eq:closure-right-end}
\end{align}
Moreover, if $x>0$, $z<1$, and $0\le p\le p_*$, replacing $z$ by
$1$ in $F(p)$ gives a number at least $1+s$.
\end{lemma}

\begin{proof}
First suppose $m\ne1$.  If $x>0$ and $z<1$, the chains imply
\begin{equation}
 w\ge x,\qquad y\ge x+s,\qquad z\ge y,\qquad z\ge w+s.
 \label{eq:closure-chain-box}
\end{equation}
Indeed $x>0$ first gives positive $y,w$; if $y$ or $w$ were $1$,
the corresponding chain to $z$ would force $z=1$.  The non-saturated
branches of the time chains then give $y\ge x+s$ and $z\ge w+s$.
The other two chains give $w\ge x$ and $z\ge y$.  Thus
\begin{equation}
 0\le x\le w\le1-s,\qquad x+s\le y\le1.
 \label{eq:closure-box-bounds}
\end{equation}
The midpoint identity
\[
 F(p_*)-1
 =\frac{(y-x)+(z-w)+m(w-x)+m(z-y)}2
\]
proves \eqref{eq:closure-midpoint} for every $m>0$, including
$m<1$.  If instead $x>0,z=1$, the first time chain gives $y\ge s$,
so
\[
 F(0)=1+y+mw\ge1+s.
\]
If $x=0,z<1$, the chain $w\to z$ gives $w\le1-s$ whenever $w>0$;
this bound also holds for $w=0$.  Hence
\[
 F(1+m)=2+m-my-w\ge2-w\ge1+s.
\]

For the last assertion write the expression after replacement as
\[
 A=1+p-px+(1-p)y+(m-p)w.
\]
If $0\le p\le1$ and $p\le p_*$, the exact identity
\begin{align*}
 A-(1+s)
  ={}&(1-p)(y-x-s)+(1+m-2p)x\\
     &+m(w-x)+p(1-s-w)
\end{align*}
is a sum of nonnegative terms by \eqref{eq:closure-box-bounds}.
In particular it covers the whole range $0\le p\le p_*$ when
$m<1$, even when $m-p<0$ in the original expression for $A$.
If $1<p\le p_*$, necessarily $m\ge2p-1>p$, and the exact identity
\begin{align*}
 A-(1+s)
  ={}&(p-1)(1-y)+(m-p)(w-x)\\
     &+(1+m-2p)x+(1-s-x)
\end{align*}
again has nonnegative terms.  This covers the remaining range for
$m>1$ and agrees with the first calculation at $p=1$.

For $m=1$, use $q$ rather than separate middle rates.  The three
relevant evaluations are
\[
 F(1)=1+z-x,\qquad
 F(0)\big|_{z=1}=1+q,\qquad
 F(2)\big|_{x=0}=3-q.
\]
The tied-clock implications $z\ge x+s$ when $x>0,z<1$,
$q\ge s$ when $x>0$, and $q\le2-s$ when $z<1$ prove
\eqref{eq:closure-midpoint}--\eqref{eq:closure-right-end}.
For $0\le p\le1=p_*$ the expression with $z$ replaced by $1$ is
$A=1+p-px+(1-p)q$, and
\[
 A-(1+s)=p(1-s-x)+(1-p)(q-s)\ge0.
\]
Here $x\le1-s$ follows from $z\ge x+s$ and $z<1$.
\end{proof}

\subsection{Contradiction and completion}

\begin{proof}[Completion of the proof of Theorem~\ref{thm:classical-growth}]
Choose a contact point satisfying the preceding lemmas, and keep its
$p,s$ fixed.  The identity \eqref{eq:closure-contact-identity} has
right side strictly below the algebraic threshold:
\begin{equation}
 2s-\beta-\gamma m<1+s,
 \label{eq:closure-strict-gap}
\end{equation}
since $s\le1$ and $\beta>\gamma>0$.  All label rates $D$ are
nonnegative.

Suppose first that $p\ge p_*$.  Since $p>0$, the point cannot lie
on the lower aperture boundary, so the minus side is admissible.
Choose there a tuple with $B\le0$.  Lemma~\ref{lem:closure-positive-x}
gives $x=x^-\ge\gamma>0$.  For this tuple $F$ is nondecreasing.
If $z<1$, then $F(p)\ge F(p_*)\ge1+s$; if $z=1$, then
$F(p)\ge F(0)\ge1+s$.  Both alternatives contradict
\eqref{eq:closure-contact-identity} and \eqref{eq:closure-strict-gap}.
This includes $p=p_*$ and the upper boundary value $p=1+m$.
At the latter value $z$ can lack a strong trace, but its coefficient
in $F(p)$ is zero and the selected tuple is sufficient.

It remains that $0\le p<p_*$.  In this range $p<1+m$, so the
point cannot be on the upper aperture boundary and the plus side is
admissible.  Choose a plus tuple with $B\ge0$, for which $F$ is
nonincreasing.  If it had $z<1$, the case $x>0$ would give
$F(p)\ge F(p_*)\ge1+s$, and the case $x=0$ would give
$F(p)\ge F(1+m)\ge1+s$.  Either is impossible.  Thus this tuple
has $z=1$.  Since the $Z$ clock is transverse for $p<1+m$, the
fixed plus trace satisfies $z^+=1$, and every plus tuple has that
same value.

If $p=0$, take the plus tuple with $x>0$ from
Lemma~\ref{lem:closure-positive-x}.  It too has $z=1$, so
$F(0)\ge1+s$, a contradiction.  This handles the lower aperture
boundary even though $X$ is then characteristic and need not have a
strong trace.

Finally suppose $0<p<p_*$.  Both sides are admissible, and
$x^->0$.  The $Z$ rate has strong traces on both sides.  If
$z^-=1$, choose a minus tuple with $B\le0$; its nondecreasing $F$
would satisfy $F(p)\ge F(0)\ge1+s$.  Therefore $z^-<1$.
For any minus tuple the jump estimate, $z^+=1$, and $D^+\ge0$ give
\[
 D^-\ge(1+m-p)(1-z^-).
\]
Consequently
\[
 F(p)+D^-\ge F(p)+(1+m-p)(1-z^-).
\]
The right side is exactly $F(p)$ with its $z$ entry replaced by $1$.
It is at least $1+s$ by Lemma~\ref{lem:closure-sign-algebra}, once
more contradicting \eqref{eq:closure-contact-identity} and
\eqref{eq:closure-strict-gap}.

The three ranges $m<1$, $m=1$, and $m>1$ have all been included in
the algebraic lemma.  At the intermediate values $p=1$ or $p=m$,
an untraced middle component has zero tangential coefficient; the
tuple argument and the nonnegative identities above remain valid.
Thus no matched-clock least-aspect extremizer can exist under
$\beta>\gamma$.  The isotropic alternative was already excluded by
Proposition~\ref{prop:isotropic-exclusion}.  The contradiction proves
$\beta\le\gamma$, completing Theorem~\ref{thm:classical-growth}.
\end{proof}
